\documentclass[11pt]{article}
\newcommand{\nombredocumento}{03 · Guía de fabricación}
% Preámbulo compartido por todos los documentos del proyecto (compilar con pdflatex).
\usepackage[utf8]{inputenc}
\usepackage[T1]{fontenc}
\usepackage{lmodern}
\usepackage{textcomp}
\usepackage[spanish,es-nodecimaldot,es-noshorthands,es-tabla]{babel}
\DeclareUnicodeCharacter{03A9}{\ensuremath{\Omega}}
\DeclareUnicodeCharacter{00B5}{\ensuremath{\mu}}
\DeclareUnicodeCharacter{03BC}{\ensuremath{\mu}}
\DeclareUnicodeCharacter{221A}{\ensuremath{\surd}}
\DeclareUnicodeCharacter{2192}{\ensuremath{\rightarrow}}
\DeclareUnicodeCharacter{2248}{\ensuremath{\approx}}
\DeclareUnicodeCharacter{2212}{\ensuremath{-}}
\DeclareUnicodeCharacter{2264}{\ensuremath{\le}}
\DeclareUnicodeCharacter{2265}{\ensuremath{\ge}}
\DeclareUnicodeCharacter{00D7}{\ensuremath{\times}}
\usepackage[a4paper,top=2.5cm,bottom=2.5cm,left=2.4cm,right=2.4cm,headheight=15pt]{geometry}
\setlength{\emergencystretch}{3em}
\usepackage{amsmath,amssymb,mathtools,bm}
\usepackage{siunitx}
\sisetup{output-decimal-marker={.},group-digits=false,per-mode=symbol,separate-uncertainty}
\usepackage{graphicx,float,booktabs,tabularx,longtable,array,multirow}
\usepackage[font=small,labelfont=bf]{caption}
\usepackage{enumitem}
\setlist{itemsep=2pt,topsep=4pt}
\usepackage[dvipsnames,table]{xcolor}
\usepackage[most]{tcolorbox}
\usepackage{tikz}
\usetikzlibrary{arrows.meta,positioning,shapes.geometric,calc,fit,backgrounds,decorations.markings}
\usepackage[european,straightvoltages]{circuitikz}
\usepackage{listings}
\usepackage{fancyhdr}
\usepackage{titlesec}
\usepackage{xurl}
\usepackage[colorlinks=true,linkcolor=Azul,urlcolor=Azul,citecolor=Azul]{hyperref}

\graphicspath{{../../simulacion/graficas/}{figuras/}}

% ---------------------------------------------------------------- colores
\definecolor{Azul}{HTML}{0D47A1}
\definecolor{Cian}{HTML}{00838F}
\definecolor{Naranja}{HTML}{E65100}
\definecolor{Rojo}{HTML}{C62828}
\definecolor{Verde}{HTML}{2E7D32}
\definecolor{Gris}{HTML}{455A64}
\definecolor{FondoCodigo}{HTML}{F5F7FA}

% ---------------------------------------------------------------- títulos
\titleformat{\section}{\Large\bfseries\color{Azul}}{\thesection}{0.8em}{}
\titleformat{\subsection}{\large\bfseries\color{Cian}}{\thesubsection}{0.8em}{}
\titleformat{\subsubsection}{\normalsize\bfseries\color{Gris}}{\thesubsubsection}{0.8em}{}

% ---------------------------------------------------------------- cabeceras
\pagestyle{fancy}
\fancyhf{}
\fancyhead[L]{\small\color{Gris}\nombredocumento}
\fancyhead[R]{\small\color{Gris}Generador de flujo axial con imanes NdFeB}
\fancyfoot[C]{\small\thepage}
\renewcommand{\headrulewidth}{0.4pt}

% ---------------------------------------------------------------- cajas
\newtcolorbox{nota}[1][Nota]{enhanced,breakable,colback=Azul!4,colframe=Azul!70,
  fonttitle=\bfseries,title={#1},left=6pt,right=6pt}
\newtcolorbox{advertencia}[1][Advertencia]{enhanced,breakable,colback=Rojo!4,colframe=Rojo!80,
  fonttitle=\bfseries,title={#1},left=6pt,right=6pt}
\newtcolorbox{resultado}[1][Resultado]{enhanced,breakable,colback=Verde!5,colframe=Verde!75,
  fonttitle=\bfseries,title={#1},left=6pt,right=6pt}
\newtcolorbox{paso}[1]{enhanced,breakable,colback=white,colframe=Cian!80,
  fonttitle=\bfseries,title={#1},left=6pt,right=6pt}

% ---------------------------------------------------------------- código
\lstdefinestyle{base}{
  basicstyle=\ttfamily\footnotesize,
  backgroundcolor=\color{FondoCodigo},
  frame=single,rulecolor=\color{Gris!40},
  breaklines=true,columns=fullflexible,keepspaces=true,
  keywordstyle=\color{Azul}\bfseries,
  commentstyle=\color{Verde}\itshape,
  stringstyle=\color{Naranja},
  showstringspaces=false,
  extendedchars=true,
  literate={á}{{\'a}}1 {é}{{\'e}}1 {í}{{\'i}}1 {ó}{{\'o}}1 {ú}{{\'u}}1 {ñ}{{\~n}}1
           {Á}{{\'A}}1 {É}{{\'E}}1 {Í}{{\'I}}1 {Ó}{{\'O}}1 {Ú}{{\'U}}1 {Ñ}{{\~N}}1
           {·}{{$\cdot$}}1 {→}{{$\rightarrow$}}1 {Ø}{{\O}}1 {Ω}{{$\Omega$}}1 {µ}{{$\mu$}}1
           {−}{{-}}1 {²}{{$^2$}}1 {√}{{$\surd$}}1 {≈}{{$\approx$}}1
}
\lstdefinestyle{python}{style=base,language=Python}
\lstdefinelanguage{JavaScript}{
  keywords={const,let,var,function,return,class,constructor,new,if,else,for,while,this,true,false,null,typeof,module,export,import},
  sensitive=true,comment=[l]{//},morecomment=[s]{/*}{*/},morestring=[b]',morestring=[b]",morestring=[b]`
}
\lstdefinestyle{js}{style=base,language=JavaScript}
\lstdefinestyle{consola}{style=base,keywordstyle=\color{black}}

% ---------------------------------------------------------------- utilidades
\newcommand{\vect}[1]{\mathbf{#1}}
\newcommand{\unitario}[1]{\mathbf{a}_{#1}}
\newcommand{\dif}{\mathrm{d}}
\DeclareUrlCommand\archivo{\urlstyle{tt}}
\makeatletter
\newcommand{\filasdesde}[1]{\@@input #1 }
\makeatother

\newcommand{\portada}[3]{%
  \begin{titlepage}
    \begin{tikzpicture}[remember picture,overlay]
      \fill[Azul] (current page.north west) rectangle ([yshift=-7cm]current page.north east);
      \node[anchor=north west,text=white,font=\sffamily\bfseries\Huge,text width=16cm,align=left]
        at ([xshift=2.4cm,yshift=-2.2cm]current page.north west) {#1};
      \node[anchor=north west,text=white!85!Azul,font=\sffamily\Large,text width=16cm,align=left]
        at ([xshift=2.4cm,yshift=-4.6cm]current page.north west) {#2};
    \end{tikzpicture}
    \vspace*{7cm}
    \begin{center}
      \begin{tikzpicture}[scale=0.9]
        \foreach \i in {0,...,7} {
          \pgfmathsetmacro{\ang}{45*\i}
          \ifodd\i \def\col{Azul!70}\def\pol{S}\else \def\col{Rojo!80}\def\pol{N}\fi
          \fill[\col] (\ang:2.2) circle (0.55);
          \node[white,font=\bfseries] at (\ang:2.2) {\pol};
        }
        \draw[Gris,line width=1pt] (0,0) circle (3.0);
        \fill[Gris!40] (0,0) circle (0.35);
      \end{tikzpicture}
    \end{center}
    \vfill
    {\large #3\par}
    \vspace{0.6cm}
    {\small\color{Gris}Basado en M.\,N.\,O. Sadiku, \emph{Elementos de Electromagnetismo}, 3.ª ed., Oxford University Press.\par
    Valores numéricos generados por el paquete \texttt{generador\_axial} (\texttt{uv run generador-axial reporte}).\par}
  \end{titlepage}
}

% Archivo generado automáticamente por `generador-axial reporte`. No editar a mano.
\newcommand{\ValGradoIman}{N38}  % Grado comercial del imán
\newcommand{\ValBr}{1.25}  % Remanencia [T]
\newcommand{\ValDiametroIman}{20}  % [mm]
\newcommand{\ValRadioIman}{10.0}  % [mm]
\newcommand{\ValEspesorIman}{5.0}  % [mm]
\newcommand{\ValAreaImanMM}{314.16}  % [mm^2]
\newcommand{\ValPolos}{8}
\newcommand{\ValParesPolos}{4}
\newcommand{\ValPasoPolarGrados}{45.0}  % [grados]
\newcommand{\ValRadioMontaje}{40.0}  % [mm]
\newcommand{\ValBobinas}{8}
\newcommand{\ValEspiras}{180}
\newcommand{\ValEspirasTotales}{1440}
\newcommand{\ValDiametroInteriorBobina}{12.0}  % [mm]
\newcommand{\ValDiametroExteriorBobina}{28.0}  % [mm]
\newcommand{\ValAlturaBobina}{6.0}  % [mm]
\newcommand{\ValEntrehierro}{2.5}  % [mm]
\newcommand{\ValEntrehierroMasAltura}{8.5}  % [mm]
\newcommand{\ValAcoplamiento}{0.85}
\newcommand{\ValCalibreAWG}{26}
\newcommand{\ValDiametroCobre}{0.404}  % [mm]
\newcommand{\ValDiametroEsmaltado}{0.455}  % [mm]
\newcommand{\ValAreaCobreMM}{0.1282}  % [mm^2]
\newcommand{\ValResistividadCobre}{1.72}  % [1e-8 ohm m]
\newcommand{\ValLongitudMediaEspira}{62.83}  % [mm]
\newcommand{\ValLongitudAlambreBobina}{11.31}  % [m]
\newcommand{\ValLongitudAlambreTotal}{90.48}  % [m]
\newcommand{\ValResistenciaBobina}{1.517}  % [ohm]
\newcommand{\ValResistenciaInterna}{12.14}  % [ohm]
\newcommand{\ValVentanaBobinaMM}{48.0}  % [mm^2]
\newcommand{\ValOcupacion}{0.78}
\newcommand{\ValRaizCaraIman}{11.1803}  % sqrt(L^2+R^2) [mm]
\newcommand{\ValTerminoCaraIman}{0.4472}  % L/sqrt(L^2+R^2)
\newcommand{\ValCampoCaraIman}{0.2795}  % [T]
\newcommand{\ValEntrehierroMasEspesor}{7.5}  % g+L [mm]
\newcommand{\ValRaizGL}{12.5000}  % sqrt((g+L)^2+R^2) [mm]
\newcommand{\ValRaizG}{10.3078}  % sqrt(g^2+R^2) [mm]
\newcommand{\ValTerminoUnoG}{0.6000}  % (g+L)/sqrt((g+L)^2+R^2)
\newcommand{\ValTerminoDosG}{0.2425}  % g/sqrt(g^2+R^2)
\newcommand{\ValCampoCaraBobina}{0.2234}  % [T]
\newcommand{\ValEntrehierroAlturaEspesor}{13.5}  % g+h+L [mm]
\newcommand{\ValRaizGHL}{16.8003}  % sqrt((g+h+L)^2+R^2) [mm]
\newcommand{\ValRaizGH}{13.1244}  % sqrt((g+h)^2+R^2) [mm]
\newcommand{\ValPrimitivaEnG}{2.1922}  % F(g) [mm]
\newcommand{\ValPrimitivaEnGH}{3.6759}  % F(g+h) [mm]
\newcommand{\ValDiferenciaPrimitivas}{1.4837}  % F(g+h)-F(g) [mm]
\newcommand{\ValCampoMedioBobina}{0.1545}  % [T]
\newcommand{\ValReduccionPorPromedio}{30.8}  % [%]
\newcommand{\ValFlujoMaximoMicro}{41.270}  % [uWb]
\newcommand{\ValAreaImanM}{3.1416}  % [1e-4 m^2]
\newcommand{\ValCaidaDiodo}{0.30}  % [V]
\newcommand{\ValCaidaPuente}{0.60}  % [V]
\newcommand{\ValCapacitanciaMicro}{220}  % [uF]
\newcommand{\ValTensionNominalCondensador}{50}  % [V]
\newcommand{\ValRpmMaxima}{1200}
\newcommand{\ValTensionDcMaxima}{29.27}  % [V]
\newcommand{\ValMargenCondensador}{41}  % [%]
\newcommand{\ValResistenciaLed}{2200}  % [ohm]
\newcommand{\ValPotenciaNominalResistenciaLed}{0.50}  % [W]
\newcommand{\ValCaidaLed}{2.0}  % [V]
\newcommand{\ValCorrienteLedRefMili}{2.21}  % [mA]
\newcommand{\ValCorrienteLedMaxMili}{12.4}  % [mA]
\newcommand{\ValCorrienteLedLimiteMili}{20}  % [mA]
\newcommand{\ValPotenciaResistenciaLedMax}{0.338}  % [W]
\newcommand{\ValRizadoRef}{251.4}  % [mV]
\newcommand{\ValRpmReferencia}{300}
\newcommand{\ValOmegaMecanicaRef}{31.416}  % [rad/s]
\newcommand{\ValFrecuenciaRef}{20.00}  % [Hz]
\newcommand{\ValPeriodoRefMs}{50.00}  % [ms]
\newcommand{\ValOmegaElectricaRef}{125.66}  % [rad/s]
\newcommand{\ValTensionPicoRef}{7.47}  % [V]
\newcommand{\ValTensionRmsRef}{5.28}  % [V]
\newcommand{\ValTensionDcRef}{6.87}  % [V]
\newcommand{\ValTensionPicoPorBobinaRef}{0.933}  % [V]
\newcommand{\ValRpmCarga}{400}
\newcommand{\ValOmegaMecanicaCarga}{41.89}  % [rad/s]
\newcommand{\ValTensionPicoCarga}{9.96}  % [V]
\newcommand{\ValTensionRmsCarga}{7.04}  % [V]
\newcommand{\ValCorrienteCarga}{0.290}  % [A]
\newcommand{\ValTensionEnCarga}{3.52}  % [V]
\newcommand{\ValPotenciaCarga}{1.021}  % [W]
\newcommand{\ValPotenciaPerdidas}{1.021}  % [W]
\newcommand{\ValPotenciaMecanica}{2.042}  % [W]
\newcommand{\ValParResistenteMili}{48.7}  % [mN m]
\newcommand{\ValEficiencia}{50}  % [%]
\newcommand{\ValBateriaTension}{12}  % [V]
\newcommand{\ValBateriaCapacidad}{7}  % [Ah]
\newcommand{\ValBateriaEnergia}{84}  % [Wh]
\newcommand{\ValBateriaResistencia}{0.05}  % [ohm]
\newcommand{\ValBateriaVacia}{11.8}  % E_oc con SoC = 0 [V]
\newcommand{\ValBateriaLlena}{12.8}  % E_oc con SoC = 1 [V]
\newcommand{\ValBateriaPendiente}{1.0}  % [V]
\newcommand{\ValSocInicial}{50}  % [%]
\newcommand{\ValSocCorte}{20}  % [%]
\newcommand{\ValSocReconexion}{25}  % [%]
\newcommand{\ValTensionOcRef}{12.30}  % [V]
\newcommand{\ValEnergiaUtilRef}{25.2}  % [Wh]
\newcommand{\ValControladorEficiencia}{90}  % [%]
\newcommand{\ValControladorArranque}{3.5}  % [V]
\newcommand{\ValControladorCorrienteMax}{3.0}  % [A]
\newcommand{\ValControladorAbsorcion}{14.4}  % [V]
\newcommand{\ValRpmMicrored}{600}
\newcommand{\ValFemRmsMicrored}{10.56}  % [V]
\newcommand{\ValTensionDcMicrored}{14.34}  % [V]
\newcommand{\ValPotenciaExtraibleMicrored}{2.297}  % [W]
\newcommand{\ValPotenciaHaciaBateria}{2.067}  % [W]
\newcommand{\ValCorrienteCargaMicroredMili}{168}  % [mA]
\newcommand{\ValPotenciaMecanicaMicrored}{4.59}  % [W]
\newcommand{\ValParMicroredMili}{73.1}  % [mN m]
\newcommand{\ValHorasRecarga}{20.8}  % [h]
\newcommand{\ValInversorTension}{110}  % [V]
\newcommand{\ValInversorPico}{155.6}  % [V]
\newcommand{\ValInversorTensionAlta}{220}  % [V]
\newcommand{\ValInversorPicoAlto}{311.1}  % [V]
\newcommand{\ValInversorFrecuencia}{60}  % [Hz]
\newcommand{\ValInversorNominal}{150}  % [W]
\newcommand{\ValInversorEficiencia}{90}  % [%]
\newcommand{\ValRelacionEspiras}{13.0}  % V_pico,AC / V_bat
\newcommand{\ValPotenciaCargaAc}{15}  % [W]
\newcommand{\ValCorrienteAcMili}{136}  % [mA]
\newcommand{\ValPotenciaEntradaInversor}{16.67}  % [W]
\newcommand{\ValCorrienteInversorBateria}{1.36}  % [A]
\newcommand{\ValPotenciaNetaBateria}{-14.60}  % [W]
\newcommand{\ValAutonomiaHoras}{1.73}  % [h]
\newcommand{\ValAutonomiaParadoHoras}{1.51}  % [h]
\newcommand{\ValRpmEquilibrio}{1704}
\newcommand{\ValPotenciaCargaReducida}{5}  % [W]
\newcommand{\ValRpmEquilibrioReducida}{984}
\newcommand{\ValInerciaRotor}{2.415}  % [1e-4 kg m^2]
\newcommand{\ValInductanciaBobinaMili}{0.528}  % [mH]
\newcommand{\ValInductanciaTotalMili}{4.22}  % [mH]
\newcommand{\ValFlujoMaximoReal}{0.827}  % phi_max normalizado
\newcommand{\ValPicoDerivadaFlujo}{0.746}
\newcommand{\ValRmsDerivadaFlujo}{0.559}
\newcommand{\ValFactorCresta}{1.335}
\newcommand{\ValFactorMedia}{0.706}
\newcommand{\ValRelacionFemRealIdeal}{0.930}
\newcommand{\ValPuntosRadiales}{32}
\newcommand{\ValMuestrasPerfil}{256}
\newcommand{\ValPasoIntegracionMili}{5}  % [ms]
\newcommand{\ValParCoulombMili}{3.0}  % [mN m]
\newcommand{\ValCoefViscosoMicro}{15.0}  % [uN m s]
\newcommand{\ValCoefVentilacionNano}{15.0}  % [nN m s^2]
\newcommand{\ValParMaximoMotriz}{0.30}  % [N m]
\newcommand{\ValVelocidadVacioMotriz}{1500}  % [RPM]
\newcommand{\ValParImpulso}{0.08}  % [N m]
\newcommand{\ValDuracionImpulso}{0.5}  % [s]
\newcommand{\ValCapacidadTermica}{100}  % [J/K]
\newcommand{\ValConveccionNatural}{8}  % [W/m^2K]
\newcommand{\ValConveccionForzada}{0.15}  % [W/m^2K por rad/s]
\newcommand{\ValCoefCobre}{0.393}  % [%/K]
\newcommand{\ValCoefRemanencia}{-0.12}  % [%/K]
\newcommand{\ValTemperaturaMaximaPla}{55}  % [C]
\newcommand{\ValBateriaRCero}{25}  % [mohm]
\newcommand{\ValBateriaRPolarizacion}{30}  % [mohm]
\newcommand{\ValBateriaTau}{30}  % [s]
\newcommand{\ValBateriaPeukert}{1.12}
\newcommand{\ValBateriaEtaCoulomb}{95}  % [%]
\newcommand{\ValBateriaEtaLlena}{60}  % [%]
\newcommand{\ValBateriaSocGas}{85}  % [%]
\newcommand{\ValBateriaAutodescarga}{3}  % [%/mes]
\newcommand{\ValBateriaSocFlotacion}{99.5}  % [%]
\newcommand{\ValBateriaFactorGas}{1500}
\newcommand{\ValBateriaPendienteGas}{50}
\newcommand{\ValControladorFlotacion}{13.8}  % [V]
\newcommand{\ValConsumoControladorMili}{8}  % [mA]
\newcommand{\ValInversorVacio}{3.6}  % [W]
\newcommand{\ValInversorCoefCuadratico}{5.81}  % [1e-4 1/W]
\newcommand{\ValInversorCorte}{10.8}  % [V]
\newcommand{\ValInversorReconexion}{12.3}  % [V]
\newcommand{\ValParReferencia}{0.07}  % [N m]
\newcommand{\ValRpmRegimenRef}{483}
\newcommand{\ValPotenciaEjeRef}{2.40}  % [W]
\newcommand{\ValPotenciaElectromagneticaRef}{2.21}  % [W]
\newcommand{\ValPerdidasCobreRef}{0.986}  % [W]
\newcommand{\ValPerdidasRozamientoRef}{0.192}  % [W]
\newcommand{\ValPerdidasFoucaultMili}{1.65}  % [mW]
\newcommand{\ValPotenciaBateriaRef}{0.96}  % [W]
\newcommand{\ValEficienciaRef}{40}  % [%]
\newcommand{\ValTensionDcRefReal}{3.50}  % [V]
\newcommand{\ValCorrienteGeneradorRef}{0.281}  % [A]
\newcommand{\ValTemperaturaRegimenRef}{27.8}  % [C]
\newcommand{\ValRpmConCargaRef}{483}
\newcommand{\ValPotenciaEntradaInversorRef}{18.7}  % [W]
\newcommand{\ValCorrienteNetaConCarga}{-1.48}  % [A]
\newcommand{\ValRpmParMaximo}{1002}
\newcommand{\ValPotenciaBateriaParMaximo}{4.35}  % [W]
\newcommand{\ValRizadoBajo}{18.8}  % [V] con 220 uF
\newcommand{\ValRizadoAlto}{0.88}  % [V] con 4700 uF
\newcommand{\ValRpmPicoImpulso}{505}
\newcommand{\ValTiempoParadaImpulso}{2.7}  % [s]
\newcommand{\ValHorasHastaAbsorcion}{4.0}  % [h]
\newcommand{\ValHorasHastaFlotacion}{8.2}  % [h]
\newcommand{\ValCorrienteFlotacionMili}{21}  % [mA]



\begin{document}

\portada{Guía de fabricación paso a paso}
        {Construcción, conexión, rectificación y ensayo del generador de flujo axial sin núcleo}
        {Documento 3 de 5 · Procedimiento práctico}

\tableofcontents
\newpage

%=====================================================================
\section{Qué se va a construir}
%=====================================================================
El prototipo es un \textbf{generador de flujo axial sin núcleo} (\emph{coreless AFPMG}): un disco giratorio
(rotor) con imanes permanentes enfrenta a un disco fijo (estator) con bobinas de aire. Al girar, el flujo
magnético que atraviesa cada bobina cambia periódicamente y, por la ley de Faraday (Documento~1), aparece una
fuerza electromotriz alterna. Un puente rectificador la convierte en tensión continua que enciende un LED y,
mediante un controlador de carga, recarga una batería de \ValBateriaTension\ V que alimenta un inversor de
\ValInversorTension\ V / \ValInversorTensionAlta\ V con un bombillo de \ValPotenciaCargaAc\ W.

\begin{table}[H]
\centering
\caption{Especificación del prototipo (valores tomados del modelo).}
\begin{tabular}{ll}
\toprule
Elemento & Especificación\\
\midrule
Rotor & Disco $\varnothing$110 × 6 mm, \ValPolos\ imanes \ValGradoIman\ $\varnothing$\ValDiametroIman\ × \ValEspesorIman\ mm a radio \ValRadioMontaje\ mm\\
Polaridad del rotor & Alternada N-S-N-S-N-S-N-S (paso polar \ValPasoPolarGrados$^\circ$)\\
Estator & Disco $\varnothing$120 × 5 mm, \ValBobinas\ bobinas a radio \ValRadioMontaje\ mm\\
Bobina & \ValEspiras\ espiras AWG~\ValCalibreAWG; $\varnothing_{int}$ \ValDiametroInteriorBobina\ mm, $\varnothing_{ext}$ \ValDiametroExteriorBobina\ mm, altura \ValAlturaBobina\ mm\\
Conexión & Las \ValBobinas\ bobinas en serie, con polaridad alternada (\ValEspirasTotales\ espiras en total)\\
Entrehierro & \ValEntrehierro\ mm entre la cara de los imanes y la cara de las bobinas\\
Resistencia interna & \ValResistenciaInterna\ $\Omega$ (a 20\,$^\circ$C)\\
Salida a \ValRpmReferencia\ RPM & $V_{pico}=\ValTensionPicoRef$ V, $V_{RMS}=\ValTensionRmsRef$ V, $V_{DC}=\ValTensionDcRef$ V, $f_e=\ValFrecuenciaRef$ Hz\\
Micro-red & Batería \ValBateriaTension\ V / \ValBateriaCapacidad\ Ah y controlador de carga a \ValControladorAbsorcion\ V\\
Inversor & \ValInversorTension\ V o \ValInversorTensionAlta\ V, \ValInversorFrecuencia\ Hz, bombillo de \ValPotenciaCargaAc\ W\\
\bottomrule
\end{tabular}
\end{table}

\begin{nota}[Ventajas del diseño sin núcleo]
Al no haber hierro en el estator no existe \emph{cogging} (par de retención): el rotor gira suave desde
velocidad cero. Tampoco hay pérdidas por histéresis ni corrientes de Foucault en el núcleo, de modo que la
tensión generada es proporcional a la velocidad (comportamiento lineal) y el modelo analítico es muy fiel.
\end{nota}

\subsection{Orden de trabajo}
\begin{center}
\begin{tikzpicture}[
  fase/.style={draw=Cian,fill=Cian!8,rounded corners,minimum width=2.5cm,minimum height=1.1cm,align=center,font=\small},
  flecha/.style={-{Stealth[length=2.5mm]},thick,Gris}]
  \node[fase] (f0) {Fase 0\\Impresión 3D};
  \node[fase,right=0.8cm of f0,yshift=0.8cm] (f1) {Fase 1\\Bobinas y estator};
  \node[fase,right=0.8cm of f0,yshift=-0.8cm] (f2) {Fase 2\\Rotor};
  \node[fase,right=0.8cm of f1,yshift=-0.8cm] (f3) {Fase 3\\Ensamble};
  \node[fase,right=0.6cm of f3] (f4) {Fase 4\\Electrónica};
  \node[fase,below=0.7cm of f4] (f5) {Fase 5\\Micro-red};
  \node[fase,left=0.6cm of f5] (f6) {Fase 6\\Pruebas};
  \draw[flecha] (f0.east) -- (f1.west);
  \draw[flecha] (f0.east) -- (f2.west);
  \draw[flecha] (f1.east) -- (f3.west);
  \draw[flecha] (f2.east) -- (f3.west);
  \draw[flecha] (f3) -- (f4);
  \draw[flecha] (f4) -- (f5);
  \draw[flecha] (f5) -- (f6);
\end{tikzpicture}
\end{center}
Las fases 1 y 2 son independientes y pueden hacerse en paralelo (por ejemplo, mientras cura el epoxi del rotor
se bobinan las bobinas).

%=====================================================================
\section{Seguridad con imanes de neodimio}
%=====================================================================
\begin{advertencia}[Lea esto antes de abrir el paquete de imanes]
\begin{enumerate}
  \item \textbf{Pellizcos}: a menos de 10 cm dos imanes \ValGradoIman\ se atraen con fuerza suficiente para
        lesionar los dedos. Manipúlelos de uno en uno y mantenga los dedos fuera de la línea de impacto.
  \item \textbf{Fragilidad}: son cerámicos sinterizados recubiertos de níquel; si chocan entre sí se astillan
        y lanzan fragmentos cortantes. Use gafas de protección.
  \item \textbf{Separación}: para separar dos imanes, deslícelos lateralmente sobre el borde de una mesa; nunca
        tire de frente.
  \item \textbf{Electrónica y marcapasos}: mantenga tarjetas, teléfonos, relojes y marcapasos a más de 30 cm.
  \item \textbf{Temperatura}: el grado N pierde magnetización por encima de unos 80\,$^\circ$C. No suelde cerca
        de los imanes ni use pistola de calor sobre el rotor.
\end{enumerate}
\end{advertencia}

\begin{advertencia}[Rotación]
A velocidades de taladro (más de 600 RPM) un imán mal pegado sale despedido. Respete el tiempo de curado del
epoxi (Fase 2), use gafas y no se coloque en el plano del rotor durante los ensayos.
\end{advertencia}

%=====================================================================
\section{Materiales y herramientas}
%=====================================================================
\subsection{Lista de materiales}
\begin{table}[H]
\centering
\caption{Lista de materiales. Los costos son orientativos.}
\small
\begin{tabularx}{\textwidth}{clXcc}
\toprule
\# & Componente & Especificación & Cant. & USD\\
\midrule
1 & Imanes NdFeB & Disco $\varnothing$\ValDiametroIman\ × \ValEspesorIman\ mm, grado \ValGradoIman\ o superior & 8 (+1) & 8–12\\
2 & Alambre magneto & Cobre esmaltado AWG~\ValCalibreAWG\ ($\varnothing$\ValDiametroCobre\ mm de cobre); se usan \ValLongitudAlambreTotal\ m $\approx$ 104 g & carrete 150 g & 5–8\\
3 & Rodamientos & 608-ZZ o 608-2RS ($\varnothing$8 × $\varnothing$22 × 7 mm) & 2 & 2–4\\
4 & Eje & Varilla roscada o eje liso $\varnothing$8 × 150 mm, 6 tuercas y 6 arandelas M8 & 1 & 2\\
5 & Discos & Impresos en PLA/PETG (STL incluidos) o acrílico de 5–6 mm & 2 & 3–6\\
6 & Epoxi & Resina bicomponente de curado lento & 1 tubo & 3\\
7 & Diodos & 1N5819 (Schottky, $V_D\approx\ValCaidaDiodo$ V) o 1N4007 ($V_D\approx0.7$ V) & 4 & 1\\
8 & Condensador & Electrolítico \ValCapacitanciaMicro\ $\mu$F × \ValTensionNominalCondensador\ V & 1 & 0.5\\
9 & Indicador & LED de alto brillo + resistencia \ValResistenciaLed\ $\Omega$, \ValPotenciaNominalResistenciaLed\ W & 1 + 1 & 0.5\\
10 & Bornes & Bornes banana o clema de 2 polos & 2 juegos & 1\\
11 & Batería & Plomo-ácido sellada (SLA/AGM) \ValBateriaTension\ V / \ValBateriaCapacidad\ Ah & 1 & 15–20\\
12 & Controlador de carga & DC-DC buck-boost con límite de corriente (CC/CV) o MPPT de 12 V, arranque $\le$\ValControladorArranque\ V, salida \ValControladorAbsorcion\ V & 1 & 6–15\\
13 & Protección DC & Portafusible aéreo + fusible de 10 A, diodo Schottky 1N5822, cable AWG~14 & 1 juego & 3\\
14 & Inversor & 12 V DC $\rightarrow$ \ValInversorTension\ V o \ValInversorTensionAlta\ V, \ValInversorFrecuencia\ Hz, onda senoidal pura, \ValInversorNominal\ W, con corte por baja tensión & 1 & 20–35\\
15 & Carga AC & Tomacorriente en caja plástica, fusible de 1 A, roseta E27 y bombillo LED de \ValPotenciaCargaAc\ W & 1 juego & 5\\
\bottomrule
\end{tabularx}
\end{table}

\begin{nota}[Por qué \ValTensionNominalCondensador\ V y \ValResistenciaLed\ $\Omega$]
A \ValRpmMaxima\ RPM el modelo predice $V_{DC}=\ValTensionDcMaxima$ V. Un condensador común de 25 V se dañaría;
uno de \ValTensionNominalCondensador\ V conserva un margen del \ValMargenCondensador\,\%. Con esa tensión, una
resistencia de 1 k$\Omega$ disiparía unos 0.74 W y el LED recibiría unos 27 mA (por encima de sus
\ValCorrienteLedLimiteMili\ mA). Con \ValResistenciaLed\ $\Omega$ la corriente máxima es
\ValCorrienteLedMaxMili\ mA y la resistencia disipa \ValPotenciaResistenciaLedMax\ W, dentro de su potencia nominal
de \ValPotenciaNominalResistenciaLed\ W. El cálculo completo está en el Documento~2.
\end{nota}

\subsection{Herramientas}
\begin{itemize}
  \item Multímetro digital con escalas de tensión AC/DC, resistencia (200 $\Omega$) y, si es posible, frecuencia.
  \item Cautín, estaño y pasta fundente; tubo termorretráctil.
  \item Lija fina (grano 400–600) o cúter para retirar el esmalte del alambre.
  \item Brújula o imán de referencia con la cara N marcada.
  \item Calibrador (pie de rey) y galgas o láminas de \ValEntrehierro\ mm para fijar el entrehierro.
  \item Impresora 3D (o taladro con broca de 8 mm y sierra de copa si usa acrílico o MDF).
  \item Opcional: osciloscopio y tacómetro óptico.
\end{itemize}

%=====================================================================
\section{Fase 0 · Impresión de las piezas}
%=====================================================================
Las piezas se generan con
\begin{lstlisting}[style=consola]
uv run generador-axial cad
\end{lstlisting}
y quedan en \archivo{cad_3d/stl/} (unidades en milímetros). Todas son sólidos cerrados (\emph{watertight}),
verificados por las pruebas automáticas, de modo que cualquier \emph{slicer} las acepta sin reparación.

\begin{table}[H]
\centering
\caption{Piezas imprimibles. Los agujeros incluyen una holgura de 0.2 mm por lado para un ajuste deslizante.}
\small
\begin{tabularx}{\textwidth}{llX}
\toprule
Archivo & Pieza & Cotas\\
\midrule
\archivo{carrete_base.stl} & Base del carrete & Núcleo $\varnothing$\ValDiametroInteriorBobina\ × \ValAlturaBobina\ mm, ala $\varnothing$32 × 2 mm, agujero $\varnothing$8.4 mm\\
\archivo{carrete_tapa.stl} & Tapa del carrete & Ala desmontable $\varnothing$32 × 2 mm, agujero $\varnothing$8.4 mm\\
\archivo{disco_rotor_8_polos.stl} & Rotor & $\varnothing$110 × 6 mm, 8 alojamientos $\varnothing$20.4 × 4 mm a radio \ValRadioMontaje\ mm, eje $\varnothing$8.4 mm\\
\archivo{disco_estator_soporte.stl} & Estator & $\varnothing$120 × 5 mm, 8 alojamientos $\varnothing$28.4 × 3 mm a radio \ValRadioMontaje\ mm, centro $\varnothing$22.4 mm\\
\bottomrule
\end{tabularx}
\end{table}

\begin{paso}{Paso 0.1 · Configurar el \emph{slicer}}
PLA o PETG, altura de capa 0.2 mm, 4 perímetros, 40\,\% de relleno, \textbf{sin soportes}: todas las piezas
tienen una cara plana completa que va contra la cama. Imprima los alojamientos hacia arriba.
\end{paso}

\begin{paso}{Paso 0.2 · Verificar las cotas}
Mida con el calibrador el diámetro de un alojamiento de imán (debe ser 20.3–20.5 mm) y pruebe un imán: debe
entrar con una leve presión del dedo. Si queda flojo, el epoxi compensará; si no entra, lije el alojamiento
o reimprima con compensación de agujeros (\emph{hole horizontal expansion}) de $+0.1$ mm.
\end{paso}

%=====================================================================
\section{Fase 1 · Bobinas y estator}
%=====================================================================

\begin{paso}{Paso 1.1 · Montar el carrete}
Atornille \archivo{carrete_base} y \archivo{carrete_tapa} con un perno M8 y su tuerca. Cubra el núcleo con una
vuelta de cinta de teflón o papel encerado: así la bobina no se adhiere y se puede extraer.
\end{paso}

\begin{paso}{Paso 1.2 · Bobinar \ValEspiras\ espiras}
Deje 6 cm de alambre libre (inicio, \textbf{I}). Enrolle \textbf{\ValEspiras\ espiras} ordenadas y apretadas,
siempre en el mismo sentido. Puede sujetar el perno en un taladro a baja velocidad y contar vueltas, o usar
un contador mecánico. Deje 6 cm al final (\textbf{F}).

La ventana del carrete mide $\frac{\ValDiametroExteriorBobina-\ValDiametroInteriorBobina}{2}\times\ValAlturaBobina
=\ValVentanaBobinaMM$ mm$^2$; con alambre esmaltado de $\varnothing$\ValDiametroEsmaltado\ mm, \ValEspiras\
espiras ocupan el \ValOcupacion\ de la ventana (máximo práctico $\approx 230$ espiras).
\end{paso}

\begin{paso}{Paso 1.3 · Fijar, extraer y marcar}
Impregne la bobina con unas gotas de cianoacrilato o barniz dieléctrico y espere a que seque. Retire la tapa y
extraiga la bobina. Marque con cinta los extremos \textbf{I} y \textbf{F}. Repita hasta tener 8 bobinas
idénticas.
\end{paso}

\begin{paso}{Paso 1.4 · Medir cada bobina}
Raspe 1.5 cm de esmalte en ambos extremos (lija o cúter, hasta ver cobre brillante) y mida la resistencia. Cada
bobina tiene $\ValLongitudAlambreBobina$ m de alambre, por lo que debe medir
\[
R_{bobina}=\frac{\rho\,\ell}{A_{Cu}}=\frac{1.72\times10^{-8}\times\ValLongitudAlambreBobina}{\ValAreaCobreMM\times10^{-6}}
\approx \ValResistenciaBobina\ \Omega .
\]
Una bobina con resistencia mucho menor tiene espiras en cortocircuito (esmalte dañado); una con resistencia
infinita está cortada. Descártela y rehágala.
\end{paso}

\begin{paso}{Paso 1.5 · Pegar las bobinas en el estator}
Pegue cada bobina con epoxi en su alojamiento, \textbf{todas con la misma cara hacia arriba} (por ejemplo, la
cara en la que empezó a bobinar). Así todas quedan enrolladas en el mismo sentido y la alternancia de polaridad
se resuelve solo con el orden de conexión del paso siguiente. Numere las bobinas de 1 a 8 en sentido horario.
\end{paso}

\begin{figure}[H]
\centering
\begin{tikzpicture}[scale=0.85]
  % Rotor
  \begin{scope}
    \fill[Gris!15] (0,0) circle (2.75);
    \draw[Gris,thick] (0,0) circle (2.75);
    \fill[white] (0,0) circle (0.2);
    \foreach \i in {0,...,7} {
      \pgfmathsetmacro{\ang}{90-45*\i}
      \ifodd\i \def\col{Azul!70}\def\pol{S}\else \def\col{Rojo!80}\def\pol{N}\fi
      \fill[\col] (\ang:2.0) circle (0.5);
      \node[white,font=\bfseries] at (\ang:2.0) {\pol};
    }
    \node[below=0.2cm,font=\small] at (0,-2.75) {Rotor (cara hacia el estator)};
  \end{scope}
  % Estator
  \begin{scope}[xshift=7.5cm]
    \fill[Gris!8] (0,0) circle (3.0);
    \draw[Gris,thick] (0,0) circle (3.0);
    \draw[Gris] (0,0) circle (0.55);
    \foreach \i in {1,...,8} {
      \pgfmathsetmacro{\ang}{90-45*(\i-1)}
      \fill[Naranja!60] (\ang:2.0) circle (0.7);
      \fill[Gris!8] (\ang:2.0) circle (0.3);
      \node[font=\scriptsize\bfseries] at (\ang:2.0) {\i};
    }
    \node[below=0.2cm,font=\small] at (0,-3.0) {Estator (bobinas numeradas)};
  \end{scope}
\end{tikzpicture}
\caption{Disposición enfrentada: la bobina $k$ queda frente al imán $k$. Cuando la bobina impar ve un polo N, la
par ve un polo S; por eso sus fem tienen signos opuestos.}
\end{figure}

\begin{paso}{Paso 1.6 · Conexión en serie con polaridad alternada}
Una las bobinas así (``F2'' significa el extremo final de la bobina 2):
\[
\text{F1}\!-\!\text{F2},\quad \text{I2}\!-\!\text{I3},\quad \text{F3}\!-\!\text{F4},\quad \text{I4}\!-\!\text{I5},\quad
\text{F5}\!-\!\text{F6},\quad \text{I6}\!-\!\text{I7},\quad \text{F7}\!-\!\text{F8}.
\]
Quedan libres \textbf{I1} e \textbf{I8}: son los terminales de salida de corriente alterna. Suelde cada unión y
aíslela con termorretráctil.
\end{paso}

\begin{figure}[H]
\centering
\begin{circuitikz}[scale=0.95,transform shape]
  \draw (0,0) node[ocirc,label=above:I1]{}
    to[L,l=$B_1$] (1.9,0)
    to[L,l=$B_2$] (3.8,0)
    to[L,l=$B_3$] (5.7,0)
    to[L,l=$B_4$] (7.6,0)
    to[L,l=$B_5$] (9.5,0)
    to[L,l=$B_6$] (11.4,0)
    to[L,l=$B_7$] (13.3,0)
    to[L,l=$B_8$] (15.2,0) node[ocirc,label=above:I8]{};
  \foreach \x/\a/\b in {1.9/F1/F2,3.8/I2/I3,5.7/F3/F4,7.6/I4/I5,9.5/F5/F6,11.4/I6/I7,13.3/F7/F8}
    \node[below,font=\scriptsize,Gris] at (\x,-0.15) {\a--\b};
  \foreach \x/\s in {0.95/{I$\to$F},2.85/{F$\to$I},4.75/{I$\to$F},6.65/{F$\to$I},8.55/{I$\to$F},10.45/{F$\to$I},12.35/{I$\to$F},14.25/{F$\to$I}}
    \node[font=\scriptsize,Cian] at (\x,-0.95) {\s};
\end{circuitikz}
\caption{Cadena serie. Bajo cada bobina se indica en qué sentido la recorre la corriente: las bobinas pares se
recorren al revés, lo que invierte el signo de su fem y compensa que vean el polo opuesto.}
\end{figure}

\begin{nota}[Por qué se alterna]
En un instante dado la fem de la bobina $k$ es $\mathcal{E}_k=(-1)^{k+1}\,N\omega_e\Phi_{max}\sin(\omega_e t)$, porque
bobinas vecinas están frente a polos opuestos. Si se conectaran todas I$\to$F, la suma sería
$\sum_k\mathcal{E}_k=0$. Recorriendo las pares al revés, cada término cambia de signo y todas se suman:
$\mathcal{E}=8N\omega_e\Phi_{max}\sin(\omega_e t)$.
\end{nota}

\begin{paso}{Paso 1.7 · Verificar resistencia y polaridad}
\begin{enumerate}
  \item Mida la resistencia entre I1 e I8: debe ser $\approx\ValResistenciaInterna\ \Omega$ ($8\times\ValResistenciaBobina$).
  \item Prueba de polaridad: conecte el multímetro en DC (escala 200 mV) entre I1 e I8. Acerque rápidamente la
        cara N de un imán a la bobina 1 y anote el signo de la desviación. Repita con la cara S sobre la bobina 2:
        el signo debe ser \textbf{el mismo}. Continúe alternando N/S sobre todas las bobinas. Si alguna da signo
        contrario, sus extremos están invertidos.
\end{enumerate}
\end{paso}

%=====================================================================
\section{Fase 2 · Rotor magnético}
%=====================================================================
\begin{paso}{Paso 2.1 · Identificar las polaridades}
Acerque cada imán a la brújula: la cara que atrae el extremo sur de la aguja es la cara \textbf{N}. Marque con
rotulador 4 imanes con la cara N hacia arriba y 4 con la cara S hacia arriba. Guárdelos separados con
espaciadores de cartón de al menos 1 cm.
\end{paso}

\begin{paso}{Paso 2.2 · Insertar los imanes}
Inserte los imanes en el disco rotor en secuencia \textbf{N-S-N-S-N-S-N-S}, de uno en uno, sujetando el disco con
la otra mano lejos de los demás imanes. El alojamiento tiene 4 mm de profundidad y el imán \ValEspesorIman\ mm de
espesor: cada imán sobresale 1 mm. Compruebe la alternancia: un imán suelto colocado sobre posiciones vecinas
debe ser atraído y repelido alternadamente.
\end{paso}

\begin{paso}{Paso 2.3 · Fijar con epoxi}
Aplique epoxi en el fondo y alrededor de cada imán. Deje curar \textbf{12 a 24 h} con el disco horizontal y sin
imanes cerca (que podrían arrastrar los del rotor mientras el epoxi está blando).
\end{paso}

%=====================================================================
\section{Fase 3 · Ensamble mecánico y entrehierro}
%=====================================================================
\begin{figure}[H]
\centering
\begin{tikzpicture}[x=0.09cm,y=0.09cm]
  % Eje
  \fill[Gris!40] (-45,-4) rectangle (50,4);
  \node[font=\scriptsize] at (-38,0) {eje $\varnothing$8};
  % Soportes con rodamientos
  \foreach \xs in {-28,32} {
    \fill[Gris!70] (\xs-3.5,-11) rectangle (\xs+3.5,-4);
    \fill[Gris!70] (\xs-3.5,4) rectangle (\xs+3.5,11);
    \fill[brown!40] (\xs-6,-70) rectangle (\xs+6,-11);
    \draw (\xs-6,-70) rectangle (\xs+6,-11);
  }
  \node[font=\scriptsize,align=center] at (-28,-76) {soporte +\\rodamiento 608};
  \node[font=\scriptsize,align=center] at (32,-76) {soporte +\\rodamiento 608};
  % Rotor
  \fill[Gris!25] (-6,4) rectangle (0,55);
  \fill[Gris!25] (-6,-55) rectangle (0,-4);
  \draw (-6,4) rectangle (0,55) (-6,-55) rectangle (0,-4);
  \fill[Rojo!80] (-4,30) rectangle (1,50);
  \fill[Azul!70] (-4,-50) rectangle (1,-30);
  \node[font=\scriptsize,anchor=east] at (-7,52) {rotor};
  % Tuercas del rotor
  \fill[Gris] (-10,-6) rectangle (-6,6);
  \fill[Gris] (0,-6) rectangle (4,6);
  % Estator con bobinas
  \fill[Gris!12] (6.5,11.2) rectangle (11.5,60);
  \fill[Gris!12] (6.5,-60) rectangle (11.5,-11.2);
  \draw (6.5,11.2) rectangle (11.5,60) (6.5,-60) rectangle (11.5,-11.2);
  \fill[Naranja!60] (3.5,26) rectangle (9.5,54);
  \fill[Naranja!60] (3.5,-54) rectangle (9.5,-26);
  \fill[Gris!12] (3.5,34) rectangle (9.5,46);
  \fill[Gris!12] (3.5,-46) rectangle (9.5,-34);
  \node[font=\scriptsize,anchor=west] at (12.5,57) {estator};
  % Fijación del estator
  \draw[thick,Gris] (11.5,58) -- (32,58) -- (32,11);
  % Entrehierro
  \draw[{Stealth[length=1.5mm]}-{Stealth[length=1.5mm]},Rojo] (1,62) -- (3.5,62);
  \draw[densely dotted,Rojo] (1,50) -- (1,63) (3.5,54) -- (3.5,63);
  \node[Rojo,font=\scriptsize,anchor=south] at (2.25,63) {$g=\ValEntrehierro$ mm};
  \node[font=\scriptsize,anchor=east] at (-5,40) {imán};
  \node[font=\scriptsize,anchor=west] at (12.5,40) {bobina};
\end{tikzpicture}
\caption{Corte del ensamble (no a escala en el sentido axial del entrehierro). El rotor gira con el eje; el
estator se fija a un soporte y no toca el eje (agujero central de 22.4 mm).}
\end{figure}

\begin{paso}{Paso 3.1 · Soportes y eje}
Monte los dos rodamientos 608 en soportes rígidos (torres en ``U'' de madera, aluminio o impresas). Pase el eje
de 8 mm y fije el rotor con una tuerca y una arandela a cada lado. Gire a mano y observe el borde del disco:
no debe cabecear más de 0.3 mm (compruébelo con una galga fija junto al borde).
\end{paso}

\begin{paso}{Paso 3.2 · Montar el estator}
Fije el estator al soporte, con las bobinas mirando hacia los imanes y sin tocar el eje. Alinee los centros:
la bobina 1 y el imán 1 deben poder quedar enfrentados exactamente.
\end{paso}

\begin{paso}{Paso 3.3 · Ajustar el entrehierro}
Inserte una galga de \ValEntrehierro\ mm entre la cara de un imán y la cara de la bobina enfrentada y acerque
el estator hasta que la galga quede apenas pinzada. Repita en tres posiciones a 120$^\circ$ para que el estator
quede paralelo. Retire las galgas y apriete.

\emph{Por qué importa}: el campo cae rápidamente con la distancia,
\[
B_z(z)=\frac{B_r}{2}\left[\frac{z+L}{\sqrt{(z+L)^2+R^2}}-\frac{z}{\sqrt{z^2+R^2}}\right],
\]
de modo que con $g=\ValEntrehierro$ mm el campo en la cara de la bobina es \ValCampoCaraBobina\ T, y con
$g=5$ mm baja a 0.162 T: duplicar el entrehierro reduce el campo en la cara de la bobina un 27\,\%, y la
tensión cae en una proporción parecida.
\end{paso}

\begin{paso}{Paso 3.4 · Prueba de giro}
Gire a mano. El rotor debe girar suave, \textbf{sin tirones} (cero \emph{cogging}) y sin rozar el estator en
ningún punto. Si roza, revise el paralelismo (paso 3.3) o el cabeceo (paso 3.1).
\end{paso}

%=====================================================================
\section{Fase 4 · Rectificación y salida}
%=====================================================================
\begin{figure}[H]
\centering
\begin{circuitikz}[scale=0.88,transform shape]
  \coordinate (A) at (0,0);
  \coordinate (T) at (2,1.6);
  \coordinate (B) at (2,-1.6);
  \coordinate (C) at (4,0);
  \draw (A) to[D*,l=$D_1$] (T);
  \draw (C) to[D*,l_=$D_2$] (T);
  \draw (B) to[D*,l=$D_3$] (A);
  \draw (B) to[D*,l_=$D_4$] (C);
  % Generador
  \draw (A) -- (-1.8,0) to[sV,l_={Generador}] (-1.8,-3.6) -- (5,-3.6) -- (5,-2.85);
  \draw (5,-2.5) node[jump crossing,scale=1.6](X){};
  \draw (C) -- (5,0) -- (X.north);
  \draw (X.south) -- (5,-2.85);
  \node[left,font=\small] at (-0.1,0.25) {I1};
  \node[above,font=\small] at (4.5,0) {I8};
  % Rieles
  \draw (T) -- (2,2.6) -- (12,2.6);
  \draw (B) -- (2,-2.5) -- (X.west);
  \draw (X.east) -- (12,-2.5);
  % Filtro e indicador
  \draw (7,2.6) to[eC,l={\ValCapacitanciaMicro\,$\mu$F / \ValTensionNominalCondensador\,V}] (7,-2.5);
  \draw (10,2.6) to[R,l={\ValResistenciaLed\,$\Omega$}] (10,0) to[leD*,l={LED}] (10,-2.5);
  \draw (12,2.6) to[short,-o] (12.6,2.6) node[right]{$+V_{DC}$};
  \draw (12,-2.5) to[short,-o] (12.6,-2.5) node[right]{GND};
  \fill (7,2.6) circle (2pt) (10,2.6) circle (2pt) (7,-2.5) circle (2pt) (10,-2.5) circle (2pt);
\end{circuitikz}
\caption{Puente rectificador de onda completa con filtro capacitivo e indicador LED. El cable del generador cruza
el riel negativo sin conectarse (salto).}
\end{figure}

\begin{paso}{Paso 4.1 · Puente rectificador}
Lista de conexiones (el cátodo es el extremo con la franja):
\begin{itemize}
  \item $D_1$: ánodo en I1, cátodo en (+). \quad $D_2$: ánodo en I8, cátodo en (+).
  \item $D_3$: ánodo en (−), cátodo en I1. \quad $D_4$: ánodo en (−), cátodo en I8.
\end{itemize}
En el semiciclo en que I1 es positivo conducen $D_1$ y $D_4$; en el otro, $D_2$ y $D_3$. Siempre hay
\textbf{dos diodos en serie}, por eso la caída total es $2V_D=\ValCaidaPuente$ V con Schottky 1N5819
(1.4 V con 1N4007).
\end{paso}

\begin{paso}{Paso 4.2 · Condensador de filtro}
Conecte el condensador de \ValCapacitanciaMicro\ $\mu$F entre (+) y (−) respetando su polaridad (la franja del
encapsulado marca el terminal negativo). Su función es sostener la tensión entre picos: con la carga del LED a
\ValRpmReferencia\ RPM el rizado es de unos \ValRizadoRef\ mV (Documento~2).
\end{paso}

\begin{paso}{Paso 4.3 · Indicador LED}
Conecte la resistencia de \ValResistenciaLed\ $\Omega$ en serie con el LED (ánodo hacia la resistencia, cátodo
hacia (−)). La corriente es $I_{LED}=(V_{DC}-\ValCaidaLed)/\ValResistenciaLed$: unos \ValCorrienteLedRefMili\ mA a
\ValRpmReferencia\ RPM (brillo tenue) y \ValCorrienteLedMaxMili\ mA a \ValRpmMaxima\ RPM.
\end{paso}

\begin{paso}{Paso 4.4 · Bornes de medida}
Deje dos bornes en la salida AC (I1, I8) y dos en la salida DC (+, −) para poder medir sin desoldar.
\end{paso}

%=====================================================================
\section{Fase 5 · Micro-red: batería, controlador e inversor}
%=====================================================================
\begin{advertencia}[Tensión peligrosa]
La salida del inversor (\ValInversorTension\ V o \ValInversorTensionAlta\ V, pico de \ValInversorPicoAlto\ V) puede
ser mortal. Haga todas las conexiones con el inversor apagado, use una caja aislada para el tomacorriente, no
toque bornes descubiertos y mida sólo con puntas de categoría CAT~II o superior. Una batería en cortocircuito
entrega cientos de amperios: coloque el fusible pegado al borne positivo antes de conectar nada más.
\end{advertencia}

\begin{paso}{Paso 5.1 · Ajustar el controlador de carga}
\begin{enumerate}
  \item Alimente la entrada con una fuente de 12 V, sin batería en la salida.
  \item Ajuste la tensión de salida (potenciómetro CV) a \ValControladorAbsorcion\ V.
  \item Ponga el amperímetro en cortocircuito en la salida y ajuste el límite de corriente (CC) a
        \ValControladorCorrienteMax\ A como máximo.
  \item Si el módulo no tiene diodo antirretorno, coloque un 1N5822 en serie con la salida (+).
\end{enumerate}
\textbf{Por qué}: el generador tiene $R_{int}=\ValResistenciaInterna\ \Omega$. Un elevador de tensión fija sin límite
de corriente pide más de lo que el generador puede dar, su entrada colapsa y deja de cargar. El modelo supone
que el controlador trabaja cerca del punto de máxima potencia (Documento~1).
\end{paso}

\begin{paso}{Paso 5.2 · Conectar la batería}
\begin{enumerate}
  \item Negativo de la batería al negativo común (salida (−) del puente y del controlador).
  \item Fusible de 10 A en el borne positivo y, después del fusible, la salida (+) del controlador.
  \item Compruebe la tensión en reposo: entre \ValBateriaVacia\ V (vacía) y \ValBateriaLlena\ V (llena).
        Al \ValSocInicial\,\% debe leer unos \ValTensionOcRef\ V.
\end{enumerate}
\end{paso}

\begin{paso}{Paso 5.3 · Conectar el inversor y la carga}
\begin{enumerate}
  \item Con el inversor en OFF, conecte su entrada a la batería después del fusible, rojo al (+) y negro al (−).
        Es normal una pequeña chispa al cargarse sus condensadores.
  \item En la salida conecte el tomacorriente con un fusible de 1 A en la fase y el bombillo de
        \ValPotenciaCargaAc\ W.
  \item Configure el corte por baja tensión del inversor, si lo permite, cerca de 12.0 V (\ValSocCorte\,\% de carga).
\end{enumerate}
\end{paso}

%=====================================================================
\section{Fase 6 · Protocolo de pruebas}
%=====================================================================
\subsection{Cómo medir la velocidad}
Con tacómetro óptico, con la cámara lenta de un teléfono o, más simple, midiendo la frecuencia de la salida AC
con el multímetro: como $f_e=P\cdot\text{RPM}/120$ con $P=\ValPolos$,
\[
\text{RPM}=\frac{120\,f_e}{P}=15\,f_e .
\]

\subsection{Tabla de ensayos}
Los valores teóricos los calcula el modelo (\texttt{uv run generador-axial reporte}) con puente Schottky y sin
carga. Complete las columnas medidas.
\begin{table}[H]
\centering
\caption{Ensayos en vacío. Métodos de giro: 1) a mano lento, 2) a mano rápido, 3) taladro baja, 4) taladro media.}
\small
\begin{tabular}{ccccccccc}
\toprule
Ensayo & RPM & $f_e$ [Hz] & $V_{pico}$ [V] & $V_{RMS}$ [V] & $V_{DC}$ [V] & $V_{RMS}$ med. & $V_{DC}$ med. & LED\\
\midrule
\filasdesde{generado/tabla_ensayos.tex}
\bottomrule
\end{tabular}
\end{table}

\begin{paso}{Paso 6.1 · Ensayo en vacío}
Retire el LED (o desconecte su resistencia) para medir en vacío. Para cada velocidad mida $V_{RMS}$ entre I1 e
I8 (multímetro en AC) y $V_{DC}$ entre (+) y (−). Anote también la frecuencia si su multímetro la mide.
\end{paso}

\begin{paso}{Paso 6.2 · Comprobar la linealidad}
Grafique $V_{RMS}$ medido frente a RPM. Debe ser una recta que pasa por el origen (es la firma de un generador sin
núcleo). Una curva que se aplana indica deslizamiento del acople del taladro o medición fuera de rango del
multímetro.
\end{paso}

\begin{paso}{Paso 6.3 · Ensayo con carga resistiva (opcional)}
Conecte una resistencia de potencia de unos 12 $\Omega$ ($\approx R_{int}$) a la salida AC. A \ValRpmCarga\ RPM el
modelo predice $V_{carga}=\ValTensionEnCarga$ V e $I=\ValCorrienteCarga$ A (valores RMS), es decir,
$P_L=\ValPotenciaCarga$ W. Notará que el rotor se frena: es el par de Lenz,
$\tau=\ValParResistenteMili$ mN$\cdot$m.
\end{paso}

\begin{paso}{Paso 6.4 · Ensayo de carga de la batería}
Amperímetro DC (escala 10 A) en serie entre el controlador y la batería. Con la batería cerca del
\ValSocInicial\,\%, el modelo predice:
\begin{center}
\begin{tabular}{cccc}
\toprule
RPM & Potencia a la batería [W] & $I_{carga}$ teórica [mA] & $I_{carga}$ medida\\
\midrule
400 & 0.92 & 75 & \\
\ValRpmMicrored & \ValPotenciaHaciaBateria & \ValCorrienteCargaMicroredMili & \\
800 & 3.68 & 299 & \\
1000 & 5.74 & 467 & \\
\bottomrule
\end{tabular}
\end{center}
Al detener el generador la corriente debe caer a cero sin invertirse (diodo antirretorno).
\end{paso}

\begin{paso}{Paso 6.5 · Ensayo del inversor}
\begin{enumerate}
  \item Generador parado, inversor en ON. En el tomacorriente (escala 600 V AC) debe medir
        \ValInversorTension–120 V (o \ValInversorTensionAlta–230 V) a \ValInversorFrecuencia\ Hz.
  \item Con el bombillo de \ValPotenciaCargaAc\ W encendido, la batería entrega unos
        \ValCorrienteInversorBateria\ A (\ValPotenciaEntradaInversor\ W).
  \item Haga girar el generador a \ValRpmMicrored\ RPM: la descarga neta baja a unos
        \ValPotenciaNetaBateria\ W, pero \textbf{no se invierte}. Con \ValPotenciaCargaAc\ W el equilibrio exigiría
        \ValRpmEquilibrio\ RPM; con un bombillo de \ValPotenciaCargaReducida\ W basta con
        $\approx$\ValRpmEquilibrioReducida\ RPM.
  \item Autonomía esperada desde el \ValSocInicial\,\% hasta el corte: \ValAutonomiaHoras\ h a
        \ValRpmMicrored\ RPM y \ValAutonomiaParadoHoras\ h con el generador parado.
\end{enumerate}
\end{paso}

\subsection{Análisis de discrepancias}
\begin{enumerate}
  \item \textbf{Diodos}: $V_{DC}=V_{pico}-2V_D$. Con 1N4007 la salida será 0.8 V menor que la de la tabla.
  \item \textbf{Resistencia interna}: con carga, $V_{carga}=\mathcal{E}\,R_L/(R_{int}+R_L)$; un multímetro no
        carga el circuito, pero el LED sí lo hace levemente.
  \item \textbf{Factor de acoplamiento}: si todas las tensiones son menores en la misma proporción, estime el
        valor real como $k_{real}=\ValAcoplamiento\cdot V_{medido}/V_{te\acute{o}rico}$ y actualícelo en
        \archivo{configuracion.py}.
  \item \textbf{Entrehierro}: un entrehierro mayor que el de diseño reduce la tensión (paso 3.3).
  \item \textbf{Temperatura}: la resistencia del cobre aumenta un 0.39\,\% por grado Celsius.
  \item \textbf{Controlador real}: si la corriente de carga es mucho menor que la de la tabla, el controlador no
        trabaja en el punto de máxima potencia; con $\eta_c$ real $=P_{bat,medida}/P_{max}$ se puede recalibrar.
\end{enumerate}

\subsection{Diagnóstico de fallas}
\begin{table}[H]
\centering
\small
\begin{tabularx}{\textwidth}{XX}
\toprule
Síntoma & Causa probable y solución\\
\midrule
Tensión casi nula & Bobinas conectadas sin alternar (se cancelan): repita la prueba de polaridad del paso 1.7.\\
Tensión un 25\,\% menor que la esperada & Una bobina invertida: aporta $-1$ en vez de $+1$ y la suma queda en $6/8$ del valor.\\
Resistencia I1–I8 infinita & Unión mal soldada o esmalte sin raspar.\\
Resistencia I1–I8 muy baja & Espiras en cortocircuito; mida bobina por bobina.\\
AC correcta pero DC nula & Diodo invertido o abierto en el puente; revise la lista del paso 4.1.\\
El rotor roza o vibra & Cabeceo del eje o estator no paralelo (pasos 3.1 y 3.3).\\
El LED no enciende a baja velocidad & Normal: necesita $V_{DC}>\ValCaidaLed$ V, es decir, más de unas 100 RPM.\\
La batería no carga & $V_{DC}<\ValControladorArranque$ V (gire más rápido), fusible abierto o controlador sin límite de corriente que colapsa la entrada.\\
El inversor se apaga solo & Corte por baja tensión: la batería llegó al \ValSocCorte\,\%. Recárguela hasta el \ValSocReconexion\,\% o más.\\
\bottomrule
\end{tabularx}
\end{table}

\end{document}
